\subsection{HIGHEST ROOT}

PROPOSITION 25. Assume that R is irreducible. Let C be a chamber of R, and let B(C) = \{ \(\alpha_{1}, \ldots, \alpha_{l}\) \} be the corresponding basis.

\begin{enumerate}
\def\labelenumi{(\roman{enumi})}
\item
  There exists a root \(\tilde{\alpha} = \sum_{i=1}^{l} n_i \alpha_i\) such that, for every root \(\sum_{i=1}^{l} p_i \alpha_i\) , we have \(n_1 \geqslant p_1, n_2 \geqslant p_2, \ldots, n_l \geqslant p_l\) . In other words, R has a largest element for the ordering defined by C.
\item
  We have \(\tilde{\alpha} \in \overline{\mathrm{C}}\) .
\item
  We have \((\tilde{\alpha}|\tilde{\alpha})\geqslant (\alpha |\alpha)\) for every root \(\alpha\)
\item
  For every positive root \(\alpha'\) not proportional to \(\tilde{\alpha}\) , we have \(n(\alpha', \tilde{\alpha}) = 0\) or 1.
\end{enumerate}

\begin{enumerate}
\def\labelenumi{\arabic{enumi})}
\item
  Let \(\alpha = \sum_{i=1}^{l} n_i \alpha_i\) , \(\beta = \sum_{i=1}^{l} p_i \alpha_i\) be two maximal roots for the ordering defined by C. We shall prove that \(\alpha = \beta\) , which will establish (i).
\item
  If \((\alpha | \alpha_i) < 0\) for some index \(i\) , it follows that either \(\alpha + \alpha_i \in \mathbb{R}\) or \(\alpha = -\alpha_i\) (Cor. of Th. 1, no. 3), and both possibilities are absurd by the maximality of \(\alpha\) . Thus \((\alpha | \alpha_i) \geqslant 0\) for all \(i\) .
\item
  If \(\alpha < 0\) , then \(\alpha < -\alpha\) , which is absurd. Thus \(n_i \geqslant 0\) for all \(i\) . Let \(J\) be the set of \(i\) such that \(n_i > 0\) , and \(J'\) the complement of \(J\) in \(\{1, 2, \ldots, l\}\) . Then \(J \neq \varnothing\) . If \(J'\) were non-empty, there would exist an \(i \in J\) and an \(i' \in J'\) such that \((\alpha_i | \alpha_{i'}) < 0\) (Cor. 5 of Prop. 20, no. 7); we would then have
\end{enumerate}

\[
\left(\alpha \mid \alpha_ {i ^ {\prime}}\right) = \sum_ {i \in \mathbf {I}} n _ {i} \left(\alpha_ {i} \mid \alpha_ {i ^ {\prime}}\right) <   0
\]

since \((\alpha_{j}|\alpha_{k})\leqslant 0\) whenever \(j\) and \(k\) are distinct, which would contradict 2). Thus \(J' = \emptyset\) and \(n_i > 0\) for all \(i\) .

\begin{enumerate}
\def\labelenumi{\arabic{enumi})}
\setcounter{enumi}{3}
\tightlist
\item
  We have \((\beta \mid \alpha_i) \geqslant 0\) for all \(i\) by 2). We cannot have \((\beta \mid \alpha_i) = 0\) for all \(i\) since \(\beta \neq 0\) . We deduce from 3) that
\end{enumerate}

\[
(\beta \mid \alpha) = \sum_ {i} n _ {i} (\beta \mid \alpha_ {i}) > 0.
\]

If \(\gamma = \alpha - \beta \in \mathbb{R}\) , either \(\alpha > \beta\) or \(\beta > \alpha\) (Th. 3, no. 6), which contradicts the maximality of \(\alpha\) and \(\beta\) . Thus \(\alpha = \beta\) (Cor. of Th. 1, no. 3).

\begin{enumerate}
\def\labelenumi{\arabic{enumi})}
\setcounter{enumi}{4}
\tightlist
\item
  By 2), \(\tilde{\alpha} \in \overline{C}\) . We shall prove that \((\alpha'|\alpha') \leqslant (\tilde{\alpha}|\tilde{\alpha})\) for every \(\alpha' \in R\) . Since \(\overline{C}\) is a fundamental domain for \(W(R)\) , we can assume that \(\alpha' \in \overline{C}\) . We have \(\tilde{\alpha} - \alpha' \geqslant 0\) , so \((\tilde{\alpha} - \alpha' | x) \geqslant 0\) for all \(x \in \overline{C}\) . In particular \((\tilde{\alpha} - \alpha' |\tilde{\alpha}) \geqslant 0\) and \((\tilde{\alpha} - \alpha' |\alpha') \geqslant 0\) , hence \((\tilde{\alpha}|\tilde{\alpha}) \geqslant (\alpha'|\tilde{\alpha}) \geqslant (\alpha'|\alpha')\) . Thus \(n(\alpha',\tilde{\alpha})\) must be equal to 0, 1 or -1 if \(\alpha'\) is not proportional to \(\tilde{\alpha}\) . If \(\alpha' \geqslant 0\) , then \((\tilde{\alpha}|\alpha') \geqslant 0\) by 2), so \(n(\alpha',\tilde{\alpha}) \geqslant 0\) and \(n(\alpha',\tilde{\alpha})\) must be either 0 or 1. Q.E.D.
\end{enumerate}

Remark. The root

\[
\tilde {\alpha} = \sum_ {i} n _ {i} \alpha_ {i}
\]

in (i) is said to be the highest root of R (with respect to C). Note that, by (i), we have \(n_{i} \geqslant 1\) for all i.

